\documentclass[10pt,a4paper]{amsart}
\usepackage{amsmath,amssymb,amsthm,graphicx}
% Retained minimal preamble for the Stacks Project chapter more-algebra.tex (More on Algebra).
% The chapter file pinned-original-source.tex carries only \input{preamble} on its first line, so no
% preamble text lives inside the chapter itself.  The Stacks Project uses one single shared preamble for
% every chapter; the snapshot stages it identically as sources/stacks-chow/preamble.tex and
% sources/stacks-etale-cohomology/preamble.tex.  The block below is copied mechanically, line for line,
% from sources/stacks-chow/preamble.tex lines 14-18 (the Xy-pic setup) and 154-217 (the native theorem
% declarations together with the mathematical macros).  Nothing was added, removed or redefined; the
% wrapper class and the amsmath/amssymb/amsthm/graphicx base packages are supplied by the assembler.
% Chapter: more-algebra.tex -- More on Algebra.
% For commutative diagrams we use Xy-pic
\usepackage[all]{xy}

% We use 2cell for 2-commutative diagrams.
\xyoption{2cell}

% Theorem environments.
%
\theoremstyle{plain}
\newtheorem{theorem}[subsection]{Theorem}
\newtheorem{proposition}[subsection]{Proposition}
\newtheorem{lemma}[subsection]{Lemma}

\theoremstyle{definition}
\newtheorem{definition}[subsection]{Definition}
\newtheorem{example}[subsection]{Example}
\newtheorem{exercise}[subsection]{Exercise}
\newtheorem{situation}[subsection]{Situation}

\theoremstyle{remark}
\newtheorem{remark}[subsection]{Remark}
\newtheorem{remarks}[subsection]{Remarks}

\numberwithin{equation}{subsection}

% Macros
%
\def\lim{\mathop{\mathrm{lim}}\nolimits}
\def\colim{\mathop{\mathrm{colim}}\nolimits}
\def\Spec{\mathop{\mathrm{Spec}}}
\def\Hom{\mathop{\mathrm{Hom}}\nolimits}
\def\Ext{\mathop{\mathrm{Ext}}\nolimits}
\def\SheafHom{\mathop{\mathcal{H}\!\mathit{om}}\nolimits}
\def\SheafExt{\mathop{\mathcal{E}\!\mathit{xt}}\nolimits}
\def\Sch{\mathit{Sch}}
\def\Mor{\mathop{\mathrm{Mor}}\nolimits}
\def\Ob{\mathop{\mathrm{Ob}}\nolimits}
\def\Sh{\mathop{\mathit{Sh}}\nolimits}
\def\NL{\mathop{N\!L}\nolimits}
\def\CH{\mathop{\mathrm{CH}}\nolimits}
\def\proetale{{pro\text{-}\acute{e}tale}}
\def\etale{{\acute{e}tale}}
\def\QCoh{\mathit{QCoh}}
\def\Ker{\mathop{\mathrm{Ker}}}
\def\Im{\mathop{\mathrm{Im}}}
\def\Coker{\mathop{\mathrm{Coker}}}
\def\Coim{\mathop{\mathrm{Coim}}}

% Boxtimes
%
\DeclareMathSymbol{\boxtimes}{\mathbin}{AMSa}{"02}

%
% Macros for moduli stacks/spaces
%
\def\QCohstack{\mathcal{QC}\!\mathit{oh}}
\def\Cohstack{\mathcal{C}\!\mathit{oh}}
\def\Spacesstack{\mathcal{S}\!\mathit{paces}}
\def\Quotfunctor{\mathrm{Quot}}
\def\Hilbfunctor{\mathrm{Hilb}}
\def\Curvesstack{\mathcal{C}\!\mathit{urves}}
\def\Polarizedstack{\mathcal{P}\!\mathit{olarized}}
\def\Complexesstack{\mathcal{C}\!\mathit{omplexes}}
% \Pic is the operator that assigns to X its picard group, usage \Pic(X)
% \Picardstack_{X/B} denotes the Picard stack of X over B
% \Picardfunctor_{X/B} denotes the Picard functor of X over B
\def\Pic{\mathop{\mathrm{Pic}}\nolimits}
\def\Picardstack{\mathcal{P}\!\mathit{ic}}
\def\Picardfunctor{\mathrm{Pic}}
\def\Deformationcategory{\mathcal{D}\!\mathit{ef}}
\usepackage{cleveref}
\crefname{originallemma}{original lemma}{original lemmas}
\Crefname{originallemma}{Original lemma}{Original lemmas}
\makeatletter
\newlabel{more-algebra-lemma-composition}{{0A8I}{source}{}{sourcepointer-1}{}}
\newlabel{more-algebra-lemma-composition@cref}{{[originallemma][0][]0A8I}{[1][0][]source}{}{}{}}
\makeatother
\crefname{originallemma}{original lemma}{original lemmas}
\Crefname{originallemma}{Original lemma}{Original lemmas}
\makeatletter
\newlabel{remark-internal-hom}{{0GWQ}{source}{}{sourcepointer-2}{}}
\newlabel{remark-internal-hom@cref}{{[originallemma][0][]0GWQ}{[1][0][]source}{}{}{}}
\makeatother
\crefname{originallemma}{original lemma}{original lemmas}
\Crefname{originallemma}{Original lemma}{Original lemmas}
\makeatletter
\newlabel{categories-remark-internal-hom-monoidal}{{0GWJ}{source}{}{sourcepointer-3}{}}
\newlabel{categories-remark-internal-hom-monoidal@cref}{{[originallemma][0][]0GWJ}{[1][0][]source}{}{}{}}
\makeatother
\title{Composition of internal Hom complexes as a map from the total tensor product}
\author{The Stacks Project Authors}
\date{stacks-project:more-algebra.tex:local-snapshot-20261012; GNU FDL 1.2 or later}
\begin{document}
\maketitle
\setcounter{subsection}{1}
\noindent\textit{Editorial note.} All delivered mathematical text is the author's own, copied from the pinned Stacks Project chapter file more-algebra.tex; no statement or proof step was rewritten, reordered or removed.\par\medskip
\noindent\textit{Editorial note.} The chapter file pinned-original-source.tex carries only \textbackslash{}input\{preamble\} on its first line, so no preamble text lives inside the chapter. The retained file sources/stacks-more-algebra/retained-preamble.tex is a mechanical, line-for-line copy of the project's single shared preamble as staged in the snapshot (Xy-pic setup plus theorem declarations and mathematical macros). Nothing was added, removed or redefined, and the excerpt serves only as preamble, never as a span coordinate.\par\medskip
\noindent\textit{Editorial note.} Both counted spans start at the \textbackslash{}begin line of their environment, so the statement and the author's proof are delivered whole; surrounding prose of the section is deliberately outside the counted spans and is carried as standard prerequisites instead.\par\medskip
\noindent\textit{Editorial note.} The citation fields (source\_id, source\_version, authors, source\_url, license.name, source\_title) are identical to those of every other card this lane has produced from this chapter, so the batch selector treats them as one source.\par\medskip
\noindent\textit{Allowed standard prerequisites (editorial).}\begin{itemize}
\item Internal Hom complex of two complexes and the total tensor product, with their sign conventions. discussed in the same section, source lines 17880-17964, before the counted statement. Supplies the complexes appearing in the statement and the sign rules used in the verification.
\item Monoidal structure on complexes and the categorical composition map between internal Homs. the author cites Categories, Remark categories-remark-internal-hom-monoidal at proof lines 17981-17984. Explains the categorical source of the map.
\item Finiteness of the sums involved in the components of an element of a total tensor product. used at proof lines 18015-18025 to show the image is well defined. Makes the componentwise assignment legitimate.
\end{itemize}\medskip
\ifdefined\hypertarget\hypertarget{sourcepointer-1}{}\fi
\noindent\textit{Source reference.} Original lemma 0A8I: focused-seeds/AGP/sources/stacks-snapshot-20261011/tags.txt, source lines 12916--12916. The lemma's own canonical tag, quoted so that readers can resolve the label inside the delivered text; it points at the counted statement itself and carries no proof content.\par\medskip
\ifdefined\hypertarget\hypertarget{sourcepointer-2}{}\fi
\noindent\textit{Source reference.} Original lemma 0GWQ: focused-seeds/AGP/sources/stacks-snapshot-20261011/tags.txt, source lines 21080--21080. Navigation pointer to a statement the counted proof refers to but does not contain; nothing load-bearing is delegated to it.\par\medskip
\ifdefined\hypertarget\hypertarget{sourcepointer-3}{}\fi
\noindent\textit{Source reference.} Original lemma 0GWJ: focused-seeds/AGP/sources/stacks-snapshot-20261011/tags.txt, source lines 21074--21074. Navigation pointer to a statement the counted proof refers to but does not contain; nothing load-bearing is delegated to it.\par\medskip
\begin{lemma}
\label{lemma-composition}
Let $R$ be a ring. Given complexes
$K^\bullet, L^\bullet, M^\bullet$
of $R$-modules there is a canonical morphism
$$
\text{Tot}\left(
\Hom^\bullet(L^\bullet, M^\bullet) \otimes_R \Hom^\bullet(K^\bullet, L^\bullet)
\right)
\longrightarrow
\Hom^\bullet(K^\bullet, M^\bullet)
$$
of complexes of $R$-modules.
\end{lemma}

\begin{proof}
Via the discussion in Remark \ref{remark-internal-hom} the existence
of such a canonical map follows from
Categories, Remark \ref{categories-remark-internal-hom-monoidal}.
We also give a direct construction.

\medskip\noindent
An element $\alpha$ of degree $n$ of the left hand side is
$$
\alpha = (\alpha^{p, q}) \in \bigoplus\nolimits_{p + q = n}
\Hom^p(L^\bullet, M^\bullet) \otimes_R \Hom^q(K^\bullet, L^\bullet)
$$
The element $\alpha^{p, q}$ is a finite sum
$\alpha^{p, q} = \sum \beta^p_i \otimes \gamma^q_i$
with
$$
\beta^p_i = (\beta^{r, s}_i)
\in \prod\nolimits_{r + s = p} \Hom_R(L^{-s}, M^r)
$$
and
$$
\gamma^q_i = (\gamma^{u, v}_i)
\in \prod\nolimits_{u + v = q} \Hom_R(K^{-v}, L^u)
$$
The map is given by sending $\alpha$ to
$\delta = (\delta^{r, v})$ with
$$
\delta^{r, v} =
\sum\nolimits_{i, s}  \beta^{r, s}_i
\circ \gamma^{-s, v}_i
\in \Hom_R(K^{-v}, M^r)
$$
For given $r + v = n$ this sum is finite as there are only finitely many
nonzero $\alpha^{p, q}$, hence only finitely many nonzero
$\beta^p_i$ and $\gamma^q_i$.
By our sign rules we have
\begin{align*}
\text{d}(\alpha^{p, q})
& =
\text{d}_{\Hom^\bullet(L^\bullet, M^\bullet)}(\alpha^{p, q})
+ (-1)^p \text{d}_{\Hom^\bullet(K^\bullet, L^\bullet)}(\alpha^{p, q}) \\
& =
\sum \Big( \text{d}_M \circ \beta^p_i \circ \gamma^q_i
- (-1)^p \beta^p_i \circ \text{d}_L \circ \gamma^q_i \Big) \\
&
\quad + (-1)^p \sum \Big( \beta^p_i \circ \text{d}_L \circ \gamma^q_i
- (-1)^q \beta^p_i \circ \gamma^q_i \circ \text{d}_K \Big) \\
& =
\sum \Big( \text{d}_M \circ \beta^p_i \circ \gamma^q_i 
-(-1)^n \beta^p_i \circ \gamma^q_i \circ \text{d}_K \Big)
\end{align*}
It follows that the rules $\alpha \mapsto \delta$ is compatible
with differentials and the lemma is proved.
\end{proof}

\end{document}
